\documentclass[10pt,letterpaper]{article}
\usepackage{cloudsuite}
\renewcommand{\DocShort}{CH. 06}
\begin{document}
\coverpage{Chapter 6: Understanding Containerization}{Cornell study notes and applied review | FOUNDATIONS}{Concepts \textbullet{} decisions \textbullet{} verification artifacts}
\sourceStatement
\tableofcontents
\newpage
\section{Learning contract}
\begin{takeaway}
Design a deployable containerized service with explicit network and image boundaries.
\end{takeaway}
\textbf{Chapter scope:} Containers, images, pods, clusters, orchestration and multi-container roles.\\
\sourceRef{6}{115}
\subsection{Mapped source sections}
\begin{itemize}
\item 6.1 Origins and Influences
\item 6.2 Fundamental Virtualization and Containerization
\item 6.3 Understanding Containers
\item 6.4 Understanding Container Images
\item 6.5 Multi-Container Types
\item 6.6 Case Study Example
\end{itemize}
\section{Cornell cue and notes matrix}
\small\setlength{\LTleft}{0pt}\setlength{\LTright}{0pt}
\begin{longtable}{>{\RaggedRight\arraybackslash}p{0.29\textwidth}>{\RaggedRight\arraybackslash}p{0.63\textwidth}}
\toprule\textbf{Cue / essential question}&\textbf{Explanatory notes and reasoning}\\\midrule\endfirsthead\toprule\textbf{Cue}&\textbf{Notes}\\\midrule\endhead
\textbf{What are the core container elements?} & Images are deployable artifacts; containers are executing instances; engines manage them; pods and host clusters organize workload placement.\\[3pt]\textcolor{petrol}{\emph{Practice:} Sketch the relationship among image, engine, container, pod and host.}\\[8pt]\midrule
\textbf{What is gained and risked by shared kernels?} & A shared kernel reduces packaging overhead but tightens the impact of kernel exposure and host configuration. Namespaces and control groups have distinct functions.\\[3pt]\textcolor{petrol}{\emph{Practice:} List five runtime limits or isolation assumptions.}\\[8pt]\midrule
\textbf{How do images support repeatability?} & Layered immutable images can be recreated from build definitions, but tags and dependencies can drift unless digests and provenance are pinned.\\[3pt]\textcolor{petrol}{\emph{Practice:} Produce an image bill of materials and digest record.}\\[8pt]\midrule
\textbf{What does orchestration handle?} & Orchestrators reconcile desired placement and availability, manage updates and connect multiple hosts. A package manager addresses a different part of the delivery workflow.\\[3pt]\textcolor{petrol}{\emph{Practice:} Separate chart/package concerns from cluster reconciliation.}\\[8pt]\midrule
\textbf{What network concerns appear?} & Host, overlay and container network scopes define addressability and traffic paths. Network policies must cover both internal and external flows.\\[3pt]\textcolor{petrol}{\emph{Practice:} Test east-west and north-south reachability.}\\[8pt]\midrule
\textbf{How do sidecar, adapter and ambassador roles differ?} & The chapter presents these as ways to divide supplementary application responsibilities across companion containers. Their utility must justify lifecycle and resource complexity.\\[3pt]\textcolor{petrol}{\emph{Practice:} Allocate proxying, translation and logging responsibilities to named containers.}\\[8pt]\midrule
\end{longtable}\normalsize
\section{Synthesis and architecture connections}
Translate the chapter vocabulary into one explicit service requirement, one boundary, one owner and one measurable test. Later mechanism and architecture chapters provide possible realizations; the requirements remain the acceptance baseline.
\begin{workbox}[title={Applied deliverable}]
Design a deployable containerized service with explicit network and image boundaries.\\[4pt]Create a one-page architecture or operations artifact with: (1) assumptions, (2) a decision and alternatives, (3) control ownership, (4) observable acceptance criteria, and (5) a failed-case test.
\end{workbox}
\section{Active recall}
\begin{enumerate}
\item What are the core container elements?
\item What is gained and risked by shared kernels?
\item How do images support repeatability?
\item What does orchestration handle?
\end{enumerate}
\subsection{Answer-check rubric}
\begin{itemize}
\item A correct response distinguishes the named concepts rather than relying on vendor names.
\item The explanation identifies an operational consequence, boundary or tradeoff.
\item At least one measurement, test or documentary artifact could demonstrate the claim.
\end{itemize}
\section{Blank Cornell capture sheet}
\textbf{Topic / source subsection:}\enspace\rule{0.65\linewidth}{0.4pt}\\[6pt]
\begin{longtable}{p{0.28\textwidth}p{0.64\textwidth}}
\toprule Cues / questions & Notes / evidence \\ \midrule
\rule{0pt}{42pt} & \\ \midrule
\rule{0pt}{42pt} & \\ \midrule
\rule{0pt}{42pt} & \\ \midrule
\rule{0pt}{42pt} & \\ \midrule
\rule{0pt}{42pt} & \\ \midrule
\end{longtable}
\textbf{Summary in three sentences:}\\[16pt]\hrulefill\\[16pt]\hrulefill\\[16pt]\hrulefill

\refs
\end{document}
